% ============================================================================
% Sprint 1, Track 5: Adiabatic / Born--Oppenheimer projection (§III.24)
%
% Three-axis row data for §IV Table~\ref{tab:projection_axes}:
%
%   Projection : Adiabatic / Born--Oppenheimer
%   Variables  : nuclear coordinate R (or, more generally, a slow
%                continuous parameter R indexing a parameter family of
%                Hamiltonians H_el(R)); mass ratio m_e / M_n as the
%                small parameter controlling the projection's validity
%   Dimension  : preserved at the projection step (V_eff(R) has the
%                same energy dimension as H_el(R)); the new variable
%                R itself carries length [L]
%   Transc.    : none AT THE PROJECTION STEP (the projection is the
%                scale-separation operation, parameter-only); the
%                resulting V_eff(R) inherits transcendentals from
%                downstream computations (Level 3 mu(R) algebraic over
%                Q(pi, sqrt(2)); Level 4 mu(rho) piecewise-smooth;
%                Level 5 composed-geometry V_eff(R) carries the
%                exchange-constant content of every prior projection
%                in the chain)
%   Layer      : Layer 1 -> Layer 2 (introduces R as a continuous
%                parameter into the previously discrete graph
%                computation; the parameterization itself is a Layer-1
%                -> Layer-2 transition, sibling of mol-frame
%                hyperspherical separation §III.12 but structurally
%                upstream of it)
%   Sibling    : §III.25 coupled-channel / adiabatic-curve solver
%                (next step in the chain: the parameterized
%                electronic problem H_el(R) gets diagonalized, and
%                the resulting U_nu(R) becomes the effective
%                potential for the slow-variable Schroedinger
%                equation; coupled-channel adds the
%                non-adiabatic P_{nu mu} = <Phi_nu | d/dR | Phi_mu>
%                coupling that the strict-adiabatic single-channel
%                approximation drops)
%
% Sprint 1 / Track 5 / 2026-05-15
% ============================================================================

\subsection{Adiabatic / Born--Oppenheimer projection}
\label{sec:proj_adiabatic_BO}
\textbf{Source:} a single combined Hamiltonian $H$ acting on two
coupled degrees of freedom whose dynamical time-scales differ by a
parametrically large ratio (the small parameter $m_e / M_n \ll 1$ for
electron--nucleus systems, $m_e / M_\text{cluster}$ for inner electrons
versus an outer cluster, etc.).  In the foundational setting this is
the molecular Hamiltonian $H = T_n(\bm{R}) + T_e(\bm{r}) +
V_{nn}(\bm{R}) + V_{ne}(\bm{r}, \bm{R}) + V_{ee}(\bm{r})$ on a tensor
product of nuclear and electronic Hilbert spaces; in the Level~3
hyperspherical setting it is the combined kinetic and potential
operator on the joint $(R, \Omega)$ space of hyperradius and
hyperangle.\\
\textbf{Target:} a \emph{two-stage} decomposition in which the fast
degree of freedom (electronic coordinates $\bm{r}$, hyperangle
$\Omega$, valence electrons) is solved at every fixed value of the
slow degree of freedom (nuclear coordinate $\bm{R}$, hyperradius $R$,
core configuration), producing a parameter family of fast-sector
eigenproblems
\begin{equation*}
  H_\text{fast}(\bm{R}) \, \Phi_\nu(\bm{r}; \bm{R})
  = E_\nu(\bm{R}) \, \Phi_\nu(\bm{r}; \bm{R}),
\end{equation*}
parameterized by the slow variable.  The slow degree of freedom then
sees the fast-sector eigenvalues $E_\nu(\bm{R})$ as an effective
potential, producing the slow-sector Schr\"odinger equation
$[T_n + E_\nu(\bm{R})] \chi(\bm{R}) = E\, \chi(\bm{R})$ on the
slow-only Hilbert space.  The full wavefunction is reconstructed as
$\Psi(\bm{r}, \bm{R}) \approx \Phi_\nu(\bm{r}; \bm{R}) \chi(\bm{R})$
(single-channel adiabatic) or as a sum
$\sum_\nu \Phi_\nu(\bm{r}; \bm{R}) \chi_\nu(\bm{R})$
(coupled-channel; see \S~\ref{sec:proj_coupled_channel}, the
sibling projection that solves the parameterized fast-sector problem
and supplies the coupling between channels).
\\
\textbf{Variables introduced:} the slow continuous parameter $\bm{R}$
itself, which enters the previously fast-sector-only graph
computation as a \emph{continuous} index over a family of
Hamiltonians.  The small parameter controlling the projection's
validity is the mass ratio $m_e / M_n$ (and, more generally, the
ratio of fast and slow time-scales).  Both quantities are content
the bare graph does not carry: the bare graph computation is
formulated at fixed nuclear configuration with no notion of a
parametric ``slow'' variable, and the mass ratio enters only when
the slow-sector kinetic operator $T_n = -\nabla_R^2 / (2 M_n)$ is
restored.\footnote{In mass-independent (Bohr-natural) atomic units
the projection still operates: it is the structural operation of
treating the nuclear configuration as a parameter rather than as a
dynamical variable, even when the mass ratio is not numerically
manifest in the units chosen.  See Paper~4 (archived) reduced-mass
convention discussion --- the adiabatic / Born--Oppenheimer
projection is what makes the electronic problem mass-independent
(at leading order) by separating the slow nuclear sector from the
fast electronic sector; the residual mass-dependence enters at the
diagonal Born--Oppenheimer correction (DBOC) and at non-adiabatic
couplings $P_{\nu\mu}$, both subleading in $m_e / M_n$.}\\
\textbf{Dimension introduced:} preserved at the projection step.  The
fast-sector effective potential $V_\text{eff}(\bm{R}) \equiv
E_\nu(\bm{R})$ has the same energy dimension $[E]$ as the original
Hamiltonian, and the new slow-sector variable $\bm{R}$ carries length
$[L]$ (already present in the projection chain via the stereographic
or molecule-frame hyperspherical projections at \S~\ref{sec:proj_stereo}
or \S~\ref{sec:proj_molframe}).  This projection does
not add a new dimension to the Layer-2 dictionary; it
\emph{reorganizes} the existing dimensional content into a slow~$\bm{R}$
$+$ fast~$\bm{r}$ split.\\
\textbf{Transcendental signature:} \textit{none at the projection step
itself}.  The Born--Oppenheimer scale-separation is parameter-only:
it indexes a family of Hamiltonians by a continuous variable without
injecting $\pi$ or any other transcendental.  Transcendentals appear
\emph{downstream} of this projection, in the solution of the
parameterized fast-sector eigenproblem
(\S~\ref{sec:proj_coupled_channel}): Level~3 hyperspherical $\mu(R)$
is
algebraic over $\mathbb{Q}(\pi, \sqrt{2})$ (CLAUDE.md~\S~12, Track~P1,
v2.0.12); Level~4 molecule-frame hyperspherical $\mu(\rho)$ is
piecewise-smooth in $\rho$, not algebraic over any closed ring
(Track~S, v2.0.13); Level~5 composed-geometry effective potentials
inherit the full transcendental content of every prior projection in
the chain.  The clean separation is: the BO projection introduces the
continuous parameter $\bm{R}$ with no transcendental cost; solving
the parameterized fast-sector eigenproblem then introduces
transcendentals according to the natural geometry of the fast sector
and the operator-order content of $H_\text{fast}(\bm{R})$.\\
\textbf{Used in:} every Level~4 molecular calculation (Paper~15;
H$_2$, HeH$^+$, every two-electron diatomic); every Level~3
hyperspherical calculation that uses the adiabatic effective potential
(Paper~13; He, Li, Be, all atomic species with one slow
hyperradius); every composed-geometry molecular calculation
(Paper~17; LiH, BeH$_2$, H$_2$O, all polyatomics, second-row
diatomics); the \emph{double}-adiabatic separation of Paper~15~\S~VII.C
where the core--valence fiber bundle introduces a third time-scale
(slowest nuclear / slow valence-hyperradius / fast valence-angular,
plus a separate slow core-hyperradius / fast core-angular on the
core fiber); all the ab initio molecular spectroscopy results of
Paper~13~\S~IX (rovibrational H$_2$); all $D_e$ benchmarks reported
across Papers~13, 15, 17; every cross-register / cross-block
calculation in the multi-focal-composition arc (May~2026) that
treats nuclear coordinates as classical parameters in
\texttt{cross\_register\_vne.py}.  This is the most heavily-used
projection in the framework's molecular pipeline; very few molecular
results do not pass through it.\\
\textbf{Empirical anchor:} H$_2$ at 96.0\% $D_e$ (Level~4
mol-frame hyperspherical + Schwartz cusp correction, Paper~15);
LiH $R_\text{eq}$ at 5.3\% error (Level~5 composed, Paper~17);
HeH$^+$ at 93.1\% $D_e$ (Paper~15~\S~VI.E); diagonal Born--Oppenheimer
correction $\Delta E_\text{DBOC} \approx +0.0015$~Ha for He (Paper~13
~\S~VI; explains the apparent variational violation $E = -2.9052$~Ha
vs.\ exact $-2.9037$~Ha at the single-channel adiabatic level by
exactly the magnitude expected from the kinetic energy cost of
angular wavefunction variation with $R$); BeH$^+$ as a transferability
test of the ab initio Phillips--Kleinman + adiabatic composition
(Paper~17).  Each of these is a quantitative test of the BO projection
in combination with the natural-geometry sector that follows it.\\
\textbf{Honest scope.}  The strict-adiabatic single-channel
approximation is the leading order in the small parameter $m_e / M_n$;
it breaks down where two electronic surfaces become near-degenerate
(conical intersections, autoionization channels, non-adiabatic
transitions) and where the kinetic energy of the slow variable is
comparable to electronic-level spacings.  The framework's response to
this breakdown is structural, not perturbative: the 2D variational
solver (Paper~13~\S~XIII, Paper~15~\S~VI.D) treats the slow and fast
variables on equal footing, eliminating the adiabatic approximation
entirely.  Track~AR (CLAUDE.md~\S~2, v2.0.23) showed that the 4-electron
LiH adiabatic solver's apparent $D_e$ violation was an artifact of
the single-channel approximation, not of the natural geometry: the 2D
variational solver gives $E_\text{min} = -7.79$~Ha (variational bound
respected, $D_e$ unbound) versus the adiabatic solver's $D_e = +0.49$~Ha
($5.3\times$ exact).  Track~BQ (CLAUDE.md~\S~2, v2.0.32) further
established that the $R_\text{eq}$ drift $+0.15$--$0.22$~bohr/$l_\text{max}$
in LiH composed geometry is \emph{not} from the adiabatic
approximation: both the adiabatic and 2D variational solvers produce
identical drift, isolating the residual to the composed-geometry PK
architecture rather than the BO projection itself.  The named extension
when adiabatic breaks down is therefore the coupled-channel /
2D-variational sibling at \S~\ref{sec:proj_coupled_channel}.\\
\textbf{Structural reading.}  Categorically distinct from
\S~\ref{sec:proj_restmass}.  The rest-mass projection is a
ring-preserving Bohr-like rescaling of a \emph{single} particle's
spectrum under $m_e \to m_\text{red}$; it modifies the foundational
energy scale (Ry) without introducing a new continuous parameter,
and the spectrum on the projected manifold is the same up to a
rational rescaling factor.  The Born--Oppenheimer projection is a
\emph{scale-separation between two coupled degrees of freedom}: it
takes a single Hilbert space $\mathcal{H}_n \otimes \mathcal{H}_e$ and
produces a parameter family $\{ \mathcal{H}_e \mid \bm{R} \in
\mathcal{M}_\text{nuc} \}$ together with a slow-sector Schr\"odinger
equation on the parameter manifold $\mathcal{M}_\text{nuc}$.  The
small parameter $m_e / M_n$ controls the projection's accuracy but is
not itself the new variable; the new variable is the slow coordinate
$\bm{R}$, which becomes a continuous index on a previously discrete
graph computation.  The two projections answer different structural
questions: rest-mass asks ``what happens to the spectrum when the
single-particle mass changes,'' while Born--Oppenheimer asks ``how
do two coupled degrees of freedom decouple at parametrically separated
time-scales.''
The two are independent and commute (rest-mass rescaling of the
foundational $m_e$ and BO-separation of slow nuclear from fast
electronic degrees of freedom can be applied in either order); they
co-appear in every reduced-mass molecular calculation, where
rest-mass rescales the electronic-sector $m_e$ to $m_\text{red}$ on
the fast side, and Born--Oppenheimer separates the slow nuclear from
the fast electronic sector across the full molecular Hamiltonian.

Structurally, this projection is the \emph{parent} of the
coupled-channel / adiabatic-curve projection at
\S~\ref{sec:proj_coupled_channel} (sibling, immediately downstream):
BO produces the parameterized electronic problem $H_\text{el}(\bm{R})$
with its parameter family of eigenproblems, and the coupled-channel
projection is the next step that solves the parameterized problem
and assembles the adiabatic potential curves
$U_\nu(\bm{R}) = E_\nu(\bm{R})$ as effective potentials for the
slow-sector dynamics, optionally including the non-adiabatic
couplings $P_{\nu\mu}(\bm{R}) = \langle \Phi_\nu | \partial / \partial
\bm{R} | \Phi_\mu \rangle$ that the strict-adiabatic single-channel
approximation drops.  The two projections are conceptually
inseparable in the molecular pipeline --- one does not invoke
Born--Oppenheimer in practice without then solving the parameterized
problem --- but they are taxonomically distinct: BO is the
\emph{scale-separation operation} (parameter-only, no transcendental
content, Layer-1 $\to$ Layer-2 transition), while coupled-channel /
adiabatic-curve is the \emph{solution operation} (introduces the
transcendental content specific to the fast-sector geometry, and
introduces the Berry-connection-like non-adiabatic coupling content
if the single-channel approximation is relaxed).

The double-adiabatic structure of Paper~15~\S~VII.C is a structural
\emph{composition} of this projection with itself at two distinct
time-scale ratios.  In LiH composed geometry the outer BO separates
nuclear (slow) from electronic (fast); within the electronic sector,
an inner BO separates core (slow on the electronic time-scale,
\emph{very} slow on the molecular time-scale) from valence (fast),
producing the core--valence fiber bundle of Paper~17~\S~II.  Iterated
composition of the BO projection is the structural mechanism by
which the framework reaches polyatomic molecules with multiple
distinct dynamical time-scales (Paper~17 LiH/BeH$_2$/H$_2$O composed
geometries; second- and third-row frozen-core diatomics, Sprint~7
v2.19.4).  The projection's own structural content does not change
under iteration; only the count of independent parameter manifolds
$\{\mathcal{M}_\text{nuc}, \mathcal{M}_\text{core}, \mathcal{M}_\text{val},
\ldots\}$ grows.
